% Options for packages loaded elsewhere
\PassOptionsToPackage{unicode}{hyperref}
\PassOptionsToPackage{hyphens}{url}
\PassOptionsToPackage{space}{xeCJK}
\documentclass[
  a4paper,
  10pt]{article}
\usepackage{xcolor}
\usepackage[margin=22mm]{geometry}
\usepackage{amsmath,amssymb}
\setcounter{secnumdepth}{-\maxdimen} % remove section numbering
\usepackage{iftex}
\ifPDFTeX
  \usepackage[T1]{fontenc}
  \usepackage[utf8]{inputenc}
  \usepackage{textcomp} % provide euro and other symbols
\else % if luatex or xetex
  \usepackage{unicode-math} % this also loads fontspec
  \defaultfontfeatures{Scale=MatchLowercase}
  \defaultfontfeatures[\rmfamily]{Ligatures=TeX,Scale=1}
\fi
\usepackage{lmodern}
\ifPDFTeX\else
  % xetex/luatex font selection
  \setmainfont[]{Helvetica}
  \ifXeTeX
    \usepackage{xeCJK}
    \setCJKmainfont[]{Hiragino Mincho ProN}
  \fi
  \ifLuaTeX
    \usepackage[]{luatexja-fontspec}
    \setmainjfont[]{Hiragino Mincho ProN}
  \fi
\fi
% Use upquote if available, for straight quotes in verbatim environments
\IfFileExists{upquote.sty}{\usepackage{upquote}}{}
\IfFileExists{microtype.sty}{% use microtype if available
  \usepackage[]{microtype}
  \UseMicrotypeSet[protrusion]{basicmath} % disable protrusion for tt fonts
}{}
\makeatletter
\@ifundefined{KOMAClassName}{% if non-KOMA class
  \IfFileExists{parskip.sty}{%
    \usepackage{parskip}
  }{% else
    \setlength{\parindent}{0pt}
    \setlength{\parskip}{6pt plus 2pt minus 1pt}}
}{% if KOMA class
  \KOMAoptions{parskip=half}}
\makeatother
\usepackage{color}
\usepackage{fancyvrb}
\newcommand{\VerbBar}{|}
\newcommand{\VERB}{\Verb[commandchars=\\\{\}]}
\DefineVerbatimEnvironment{Highlighting}{Verbatim}{commandchars=\\\{\}}
% Add ',fontsize=\small' for more characters per line
\newenvironment{Shaded}{}{}
\newcommand{\AlertTok}[1]{\textcolor[rgb]{1.00,0.00,0.00}{\textbf{#1}}}
\newcommand{\AnnotationTok}[1]{\textcolor[rgb]{0.38,0.63,0.69}{\textbf{\textit{#1}}}}
\newcommand{\AttributeTok}[1]{\textcolor[rgb]{0.49,0.56,0.16}{#1}}
\newcommand{\BaseNTok}[1]{\textcolor[rgb]{0.25,0.63,0.44}{#1}}
\newcommand{\BuiltInTok}[1]{\textcolor[rgb]{0.00,0.50,0.00}{#1}}
\newcommand{\CharTok}[1]{\textcolor[rgb]{0.25,0.44,0.63}{#1}}
\newcommand{\CommentTok}[1]{\textcolor[rgb]{0.38,0.63,0.69}{\textit{#1}}}
\newcommand{\CommentVarTok}[1]{\textcolor[rgb]{0.38,0.63,0.69}{\textbf{\textit{#1}}}}
\newcommand{\ConstantTok}[1]{\textcolor[rgb]{0.53,0.00,0.00}{#1}}
\newcommand{\ControlFlowTok}[1]{\textcolor[rgb]{0.00,0.44,0.13}{\textbf{#1}}}
\newcommand{\DataTypeTok}[1]{\textcolor[rgb]{0.56,0.13,0.00}{#1}}
\newcommand{\DecValTok}[1]{\textcolor[rgb]{0.25,0.63,0.44}{#1}}
\newcommand{\DocumentationTok}[1]{\textcolor[rgb]{0.73,0.13,0.13}{\textit{#1}}}
\newcommand{\ErrorTok}[1]{\textcolor[rgb]{1.00,0.00,0.00}{\textbf{#1}}}
\newcommand{\ExtensionTok}[1]{#1}
\newcommand{\FloatTok}[1]{\textcolor[rgb]{0.25,0.63,0.44}{#1}}
\newcommand{\FunctionTok}[1]{\textcolor[rgb]{0.02,0.16,0.49}{#1}}
\newcommand{\ImportTok}[1]{\textcolor[rgb]{0.00,0.50,0.00}{\textbf{#1}}}
\newcommand{\InformationTok}[1]{\textcolor[rgb]{0.38,0.63,0.69}{\textbf{\textit{#1}}}}
\newcommand{\KeywordTok}[1]{\textcolor[rgb]{0.00,0.44,0.13}{\textbf{#1}}}
\newcommand{\NormalTok}[1]{#1}
\newcommand{\OperatorTok}[1]{\textcolor[rgb]{0.40,0.40,0.40}{#1}}
\newcommand{\OtherTok}[1]{\textcolor[rgb]{0.00,0.44,0.13}{#1}}
\newcommand{\PreprocessorTok}[1]{\textcolor[rgb]{0.74,0.48,0.00}{#1}}
\newcommand{\RegionMarkerTok}[1]{#1}
\newcommand{\SpecialCharTok}[1]{\textcolor[rgb]{0.25,0.44,0.63}{#1}}
\newcommand{\SpecialStringTok}[1]{\textcolor[rgb]{0.73,0.40,0.53}{#1}}
\newcommand{\StringTok}[1]{\textcolor[rgb]{0.25,0.44,0.63}{#1}}
\newcommand{\VariableTok}[1]{\textcolor[rgb]{0.10,0.09,0.49}{#1}}
\newcommand{\VerbatimStringTok}[1]{\textcolor[rgb]{0.25,0.44,0.63}{#1}}
\newcommand{\WarningTok}[1]{\textcolor[rgb]{0.38,0.63,0.69}{\textbf{\textit{#1}}}}
\usepackage{longtable,booktabs,array}
\newcounter{none} % for unnumbered tables
\usepackage{calc} % for calculating minipage widths
% Correct order of tables after \paragraph or \subparagraph
\usepackage{etoolbox}
\makeatletter
\patchcmd\longtable{\par}{\if@noskipsec\mbox{}\fi\par}{}{}
\makeatother
% Allow footnotes in longtable head/foot
\IfFileExists{footnotehyper.sty}{\usepackage{footnotehyper}}{\usepackage{footnote}}
\makesavenoteenv{longtable}
\setlength{\emergencystretch}{3em} % prevent overfull lines
\providecommand{\tightlist}{%
  \setlength{\itemsep}{0pt}\setlength{\parskip}{0pt}}
\usepackage{bookmark}
\IfFileExists{xurl.sty}{\usepackage{xurl}}{} % add URL line breaks if available
\urlstyle{same}
\hypersetup{
  hidelinks,
  pdfcreator={LaTeX via pandoc}}

\author{}
\date{}

\begin{document}

\section{反復二チャネル交換系における離散Born型重みの出現}\label{ux53cdux5fa9ux4e8cux30c1ux30e3ux30cdux30ebux4ea4ux63dbux7cfbux306bux304aux3051ux308bux96e2ux6563bornux578bux91cdux307fux306eux51faux73fe}

\subsection{波束局在性移乗・準安定二状態・観測選択から得られた有限位数回帰則}\label{ux6ce2ux675fux5c40ux5728ux6027ux79fbux4e57ux6e96ux5b89ux5b9aux4e8cux72b6ux614bux89b3ux6e2cux9078ux629eux304bux3089ux5f97ux3089ux308cux305fux6709ux9650ux4f4dux6570ux56deux5e30ux5247}

\textbf{版:} 日本語完成稿 v1 \textbf{日付:} 2026年7月18日 \textbf{著者:}
木原 範昭 \textbf{ORCID:} 0009-0004-6753-4020 \textbf{Version DOI:}
10.5281/zenodo.21422471 \textbf{Concept DOI:} 10.5281/zenodo.21422470
\textbf{位置づけ:}
波の情報読出しシリーズ・有限位数回帰とBorn型重みの独立報告

\begin{center}\rule{0.5\linewidth}{0.5pt}\end{center}

\subsection{要旨}\label{ux8981ux65e8}

本研究は、Born則またはその二状態表示である \(\cos^2\)
則を再現する目的で開始されたものではない。出発点は、二つの異なる数値実験であった。第一は、低局在波と高次倍音を含む局在波の反復交換散乱において、局在性および実効倍音構造が二チャネル間を移乗し、波束の収縮または局在化に似た挙動が現れるかを調べるSystem
Aである。第二は、A/B二状態、その中間にある灰色準安定状態、弱い読出しによる保持、および強い観測によるA/B選択を調べるSystem
Bである。

両モデルは異なる物理的問いと異なる観測関数から構成されたが、後の解析により、正規化および観測作用を加える前の線形混合部分に、同じ二チャネル交換散乱核

\[
U_R=
\begin{pmatrix}
r&t\\
t&r
\end{pmatrix}
\]

を共有することが判明した。System
Aではこの線形混合後にチャネル別正規化が入り、System
Bの強いD観測では別のバックアクション写像が入るため、両モデルの全更新写像が同一であるとは置かない。System
Bにおける交換係数 \(R\)
の全域掃引では、特定の値に極めて鋭いピークが現れた。その原因を固有値解析した結果、ピークは反対称固有値が有限回で1へ戻る厳密再帰

\[
U_R^n=I
\]

によって生成されることが分かった。

本稿では、この結果を三角関数による散乱振幅のパラメータ表示へ依存させず、交換対称性から定まる対称・反対称射影

\[
P_s=\frac12
\begin{pmatrix}
1&1\\
1&1
\end{pmatrix},
\qquad
P_a=\frac12
\begin{pmatrix}
1&-1\\
-1&1
\end{pmatrix}
\]

を用いて導出する。全体位相を固定して対称固有値を1とした交換対称ユニタリ作用素は

\[
U=P_s+\zeta P_a,
\qquad |\zeta|=1
\]

と書ける。有限位数条件が \(\zeta=e^{-2\pi im/n}\)
を選ぶと、A/B基底における対角・非対角振幅は

\[
r=\frac{1+\zeta}{2},
\qquad
t=\frac{1-\zeta}{2}
\]

となる。したがって、この有限位数条件を満たす交換重みは、互いに素な整数
\(m,n\) に対して

\[
\boxed{
R_{n,m}=\cos^2\left(\frac{\pi m}{n}\right)
}
\]

で与えられる。

この式は、Born型二乗則を探索条件または評価関数として入力して得たものではない。また、上の射影作用素表示により、三角関数形を独立な仮定として置かなくても導出できる。波束局在性移乗、準安定二状態、弱読出しおよび強観測選択を模擬する反復交換系を先に構成し、数値ピークを発見した後、その閉軌道条件を解析した結果として得られた。したがって本研究が示すのは、Born則そのものの完全導出ではなく、量子測定に関連する現象と類似した挙動を持つ閉鎖二チャネル波動系において、有限回帰可能な交換重みが離散的なBorn型
\(\cos^2\) 系列として内生的に現れることである。

現行コードでは、ユニタリ散乱振幅を \(\sin\theta\) および \(\cos\theta\)
で表す実装を用いている。しかし、この表示は交換対称射影 \(P_s,P_a\)
と固有値 \(1,\zeta\)
から回収できる座標表示であり、本稿の中心定理に必要な独立仮定ではない。どの角度が完全再帰として選ばれるかは事前に与えておらず、有限位数条件が

\[
\phi_{n,m}=\frac{\pi m}{n}
\]

という離散Born角系列を選択する。実装で用いた散乱角は
\(\theta=\pi/2-\phi_{n,m}\)
である。この結果は、交換対称な二チャネル射影構造と、閉包条件によって選択される位相系列を区別し、後者が前者の許容重みを離散化することを示す。

本稿は、微細構造定数、E8または特定の物理定数との数値対応を主題としない。中心課題は、収縮に見える局在性再配分、準安定重ね合わせ、観測選択、有限位数再帰およびBorn型二乗重みが、共通の無摂動交換核
\(U\)
と、それに続くモデル固有の正規化・読出し・選択写像によって接続されたという数理構造を明確に報告することである。

\textbf{キーワード:}
Born則、二チャネル交換、有限位数再帰、波束局在性、準安定状態、観測選択、量子測定、ユニタリ作用素、離散位相

\begin{center}\rule{0.5\linewidth}{0.5pt}\end{center}

\section{1.
研究背景と目的}\label{1-ux7814ux7a76ux80ccux666fux3068ux76eeux7684}

\subsection{1.1
Born則と二状態表示}\label{11-bornux5247ux3068ux4e8cux72b6ux614bux8868ux793a}

標準量子力学において、状態 \(|\psi\rangle\) を観測状態 \(|A\rangle\)
へ射影したときの確率は、Born則

\[
P(A)=|\langle A|\psi\rangle|^2
\]

で与えられる{[}5,6{]}。

二状態系で

\section{\$\$
\textbar\textbackslash psi\textbackslash rangle}\label{-psirangle}

\textbackslash cos\textbackslash phi,\textbar A\textbackslash rangle +
e\^{}\{i\textbackslash chi\}\textbackslash sin\textbackslash phi,\textbar B\textbackslash rangle
\$\$

と書けば、直交する二状態の観測確率は

\[
P(A)=\cos^2\phi,
\qquad
P(B)=\sin^2\phi
\]

となる。

Born則は量子理論と観測結果を接続する基本則である。一方、本稿の出発点はBorn則の導出ではなかった。本研究で先に構築された数値系は、波束の局在性移乗および準安定二状態の観測選択を調べるためのものであり、\(\cos^2\)
則を目標関数として含んでいなかった。

\subsection{1.2
本研究の問い}\label{12-ux672cux7814ux7a76ux306eux554fux3044}

本研究が扱う問いは次である。

\begin{quote}
波束収縮に似た局在性再配分、A/B二状態の準安定混合、弱い読出しによる保持、および強い観測による選択を模擬する反復交換系において、Born型二乗重みは外部から確率公理として入力しなくても、交換対称性、ユニタリ性および有限位数閉包から現れ得るか。
\end{quote}

本稿は、一般的なBorn則の一意性証明または単一試行確率の完全な力学的導出を目的としない。具体的な目的は、既に構築されていた二チャネル交換モデルに対し、交換対称なユニタリ作用素を対称・反対称部分空間へ射影分解し、その有限位数条件から離散的なBorn型重み

\[
R_{n,m}=\cos^2\left(\frac{\pi m}{n}\right)
\]

が内生的に現れることを示し、その意味と限界を明確化することである。

\subsection{1.3
発見の非誘導性}\label{13-ux767aux898bux306eux975eux8a98ux5c0eux6027}

本結果の解釈では、発見の時間順序が重要である。

\begin{enumerate}
\def\labelenumi{\arabic{enumi}.}
\tightlist
\item
  波束の局在性と倍音構造が二チャネル間を移乗するSystem
  Aを構築した{[}1,2{]}。
\item
  A/B二状態、灰色準安定状態、弱読出しおよび強観測選択を扱うSystem
  Bを構築した{[}3{]}。
\item
  System Aでは局在性移乗を局所掃引し、System Bでは共通の交換散乱係数
  \(R\) による全域掃引を行った{[}4{]}。
\item
  System Bの全域掃引結果に極めて鋭いピークを観測した。
\item
  局所精密掃引および多倍長計算を行った。
\item
  作用素固有値を解析し、ピークが有限位数再帰であることを特定した{[}4{]}。
\item
  有限位数条件から、交換重みが \(\cos^2(\pi m/n)\)
  となることを認識した。
\item
  その後、この式がBorn則の二状態表示と同型であることを物理的に再解釈した。
\end{enumerate}

したがって、本研究はBorn型 \(\cos^2\)
則をコードへ埋め込み、その出力を再発見したものではない。正確には、ユニタリ二チャネル散乱の振幅構造は初期モデルに含まれていたが、どの位相角および交換重みが完全再帰として選択されるかは入力されていなかった。さらに本稿では、その三角関数表示自体を中心仮定から外し、交換対称性とユニタリ性から

\[
U=P_s+\zeta P_a,
\qquad |\zeta|=1
\]

を導く。この表示に有限位数条件を課すことで、\(\cos^2/\sin^2\)
重みは座標表示の選び方ではなく、対称・反対称固有位相の干渉から必然的に得られる。

\begin{center}\rule{0.5\linewidth}{0.5pt}\end{center}

\section{2.
二つの出発モデル}\label{2-ux4e8cux3064ux306eux51faux767aux30e2ux30c7ux30eb}

\subsection{2.1 System
A：波束局在性移乗}\label{21-system-aux6ce2ux675fux5c40ux5728ux6027ux79fbux4e57}

System
Aでは、A側を低局在波、B側を高次倍音を含む局在波として、同じ交換散乱を反復する{[}1,2{]}。

状態を波束ベクトル \(A_j,B_j\) とし、線形交換部分を

\section{\$\$ \textbackslash begin\{pmatrix\} \textbackslash widetilde
A\_\{j+1\}\textbackslash{} \textbackslash widetilde B\_\{j+1\}
\textbackslash end\{pmatrix\}}\label{-beginpmatrix-widetilde-a_j1-widetilde-b_j1-endpmatrix}

U\_R \textbackslash begin\{pmatrix\} A\_j\textbackslash{} B\_j
\textbackslash end\{pmatrix\} \$\$

と表す。

ただし、実装されたSystem
Aの一回更新は、線形混合だけではない。\(\mathcal N(x):=x/\|x\|\)
を各波束チャネルの正規化とすると、実際の更新は

\section{\$\$ \textbackslash boxed\{ F\_R(A,B)}\label{-boxed-f_rab}

\textbackslash left( \textbackslash mathcal N(rA+tB),
\textbackslash mathcal N(tA+rB) \textbackslash right) \} \$\$

である。したがって \(F_R\)
は一般には非線形であり、本稿の有限位数定理が直接適用されるのは、正規化前の線形分子

\[
(A,B)\longmapsto(rA+tB,\ tA+rB)
\]

である。System
Aで観測された局在性移乗は、この共通線形核とチャネル別正規化を合成した全写像の現象として扱う。

観測量は、代表的には

\[
L,
\qquad
N_{\mathrm{eff}},
\qquad
B_{\mathrm{to}A}
\]

である。

ここで、

\begin{itemize}
\tightlist
\item
  \(L\) は波束の局在性、
\item
  \(N_{\mathrm{eff}}\) は実効倍音次数、
\item
  \(B_{\mathrm{to}A}\) はB側初期倍音構造がA側へ移った度合い、
\end{itemize}

を表す。

この系では、中間的な交換係数において局在性および倍音構造が二チャネル間を周期的に移乗する。ある時点の一方のチャネルだけを観測すれば、広がっていた波束が局在波へ変化したように見える。この挙動は、全体系のノルム消失を伴う非ユニタリ収縮ではなく、交換干渉による局在性の再配分として構成された。

\subsection{2.2 System
B：白猫・黒猫・灰色猫準安定界面}\label{22-system-bux767dux732bux9ed2ux732bux7070ux8272ux732bux6e96ux5b89ux5b9aux754cux9762}

System Bでは、二つの複素振幅 \(a,b\) から

\[
p_A=|a|^2,
\qquad
p_B=|b|^2,
\qquad
S=p_A-p_B
\]

を定義した{[}3{]}。

状態の読み方は次である。

\begin{itemize}
\tightlist
\item
  \(S\approx+1\)：A優勢状態、白猫状態。
\item
  \(S\approx-1\)：B優勢状態、黒猫状態。
\item
  \(S\approx0\)：A/Bが均衡した灰色状態。
\end{itemize}

灰色状態のうち、小振幅で周期的に揺らぎ、弱いC読出しでは保持されるが、強いD観測ではAまたはBへ選択される状態を灰色準安定相とした。

System Bの無摂動A/B交換は \(U_R\) に従う。実装上のペア正規化は

\section{\$\$ \textbackslash mathcal
N\_\{AB\}(a,b)}\label{-mathcal-n_abab}

\textbackslash frac\{(a,b)\}\{\textbackslash sqrt\{\textbar a\textbar\^{}2+\textbar b\textbar\^{}2\}\}
\$\$

であり、厳密なユニタリ発展では恒等的である。一方、C弱読出しおよびD強観測は、この無摂動交換の後に置かれた別の読出し・反作用写像である。とくにD実装は配分差
\(S\) を

\section{\$\$ S\_\{j+1\}}\label{-s_j1}

S\_j+g\_D S\_\{D,j\}(1-S\_j\^{}2) \$\$

と更新し、A/Bいずれかへ押し出す非線形反作用を明示的に持つ。したがって強観測選択は
\(U_R\) 単独の吸引ではなく、\(U_R\) とD反作用の合成によって生じる。

このモデルは、量子力学の測定問題を直接解決するものとして構成されたのではない。目的は、

\begin{itemize}
\tightlist
\item
  二状態の中間にある準安定状態、
\item
  状態を大きく壊さない弱い読出し、
\item
  状態を一方へ変化させる強い観測、
\end{itemize}

を同じ反復交換系の中で区別できるかを調べることであった。

\subsection{2.3
二つのモデルに共通する散乱核}\label{23-ux4e8cux3064ux306eux30e2ux30c7ux30ebux306bux5171ux901aux3059ux308bux6563ux4e71ux6838}

後の解析により、System AとSystem
Bは異なる状態空間、正規化および観測関数を持ちながら、正規化・読出し・反作用を加える前の線形混合部分に同じ交換散乱核を用いていることが分かった。

\[
\boxed{
U_R=
\begin{pmatrix}
r&t\\
t&r
\end{pmatrix}
}
\]

全更新の構造は概念的に

\[
\text{System A}:\quad F_R=\mathcal N_A\circ U_R,
\]

\[
\text{System B・無摂動}:\quad U_R,
\]

\[
\text{System B・C/D}:\quad
F_C=\mathcal C\circ U_R,
\qquad
F_D=\mathcal D\circ U_R
\]

と分解できる。ここで \(\mathcal N_A\) はチャネル別正規化、\(\mathcal C\)
は弱読出し・弱反作用、\(\mathcal D\) は強い選択反作用である。System
Aは波束ベクトルの局在性および内部倍音構造を読み、System
Bは二成分複素振幅のA/B配分差および準安定性を読む。

したがって、両モデルに共通して現れる特定の \(R\)
の候補位置は交換作用素固有の構造に由来し、そのピークの見え方、局在性移乗およびA/B選択は、後段写像と観測関数によって決まる。この「根の位置を決める共通核」と「現象を可視化・選択するモデル固有写像」の分離が、本稿の実装上の中心構造である。

\begin{center}\rule{0.5\linewidth}{0.5pt}\end{center}

\section{3.
交換散乱作用素}\label{3-ux4ea4ux63dbux6563ux4e71ux4f5cux7528ux7d20}

\subsection{3.1
数値実装で用いた散乱係数座標}\label{31-ux6570ux5024ux5b9fux88c5ux3067ux7528ux3044ux305fux6563ux4e71ux4fc2ux6570ux5ea7ux6a19}

数値実装では、反射重み \(R\in[0,1]\) に対し、

\[
\theta(R):=\arcsin\sqrt R
\]

と定義する。

透過振幅 \(t\) および反射振幅 \(r\) を

\[
t=e^{i\theta}\cos\theta,
\qquad
r=-ie^{i\theta}\sin\theta
\]

と置く。

このとき、

\[
|t|^2=\cos^2\theta=1-R,
\]

\[
|r|^2=\sin^2\theta=R
\]

である。

さらに、

\[
|r|^2+|t|^2=1,
\qquad
r^*t+t^*r=0
\]

を満たすため、

\[
U_R^\dagger U_R=I
\]

となる。

したがって固定された \(R\)
に対する線形交換作用素はユニタリである。この三角関数表示は実装で用いた便利な座標であるが、\(\cos^2/\sin^2\)
系列を得るための独立な公理ではない。第4節および第5節では、交換対称なユニタリ作用素の射影分解だけから同じ係数表示を再構成する。

\subsection{3.2
状態ノルムと局在性}\label{32-ux72b6ux614bux30ceux30ebux30e0ux3068ux5c40ux5728ux6027}

線形核 \(U_R\) のユニタリ性により、二チャネル全体のノルムは保存される。

\section{\$\$ \textbackslash left\textbar{} U\_R
\textbackslash begin\{pmatrix\} A\textbackslash B
\textbackslash end\{pmatrix\}
\textbackslash right\textbar\^{}2}\label{-left-u_r-beginpmatrix-ab-endpmatrix-right2}

\textbackslash left\textbar{} \textbackslash begin\{pmatrix\}
A\textbackslash B \textbackslash end\{pmatrix\}
\textbackslash right\textbar\^{}2 \$\$

したがって、正規化前の線形核において一方の波束が局在化したように見える場合でも、それは全ノルムの消失または一点への散逸的吸引を意味しない。局在性、倍音構造またはチャネル配分が、保存された全状態の内部で再配分されている。System
Aの全写像にはチャネル別正規化が含まれるため、この全ノルム保存式をそのまま全更新へ移すことはできず、局在性移乗は
\(F_R=\mathcal N_A\circ U_R\) として評価する。

この区別は、波束収縮との類似を議論する際に重要である。本稿でいう「収縮に見える挙動」は、標準量子力学における測定後状態更新と同一であると主張するものではなく、ユニタリ交換核、チャネル別正規化および部分観測の合成によって局在性集中が現れる数理的類似を指す。

\begin{center}\rule{0.5\linewidth}{0.5pt}\end{center}

\section{4.
対称・反対称固有モード}\label{4-ux5bfeux79f0ux53cdux5bfeux79f0ux56faux6709ux30e2ux30fcux30c9}

\subsection{4.1
固有モードへの分解}\label{41-ux56faux6709ux30e2ux30fcux30c9ux3078ux306eux5206ux89e3}

二チャネル状態を

\[
X_j:=\frac{A_j+B_j}{\sqrt2},
\qquad
Y_j:=\frac{A_j-B_j}{\sqrt2}
\]

へ分解する。

\(X_j\) はチャネル交換に対して対称な成分、\(Y_j\) は反対称な成分である。

交換作用素の固有値は

\[
\lambda_s=r+t,
\qquad
\lambda_a=r-t
\]

である。

定義した \(r,t\) を代入すると、

\[
\lambda_s=1,
\]

\[
\lambda_a=-e^{2i\theta}
\]

を得る。

したがって、

\[
X_j=X_0,
\qquad
Y_j=\lambda_a^jY_0
\]

である。

\subsection{4.2
射影作用素表示}\label{42-ux5c04ux5f71ux4f5cux7528ux7d20ux8868ux793a}

チャネル交換作用素を

\[
X=
\begin{pmatrix}
0&1\\
1&0
\end{pmatrix}
=P_s-P_a
\]

とする。交換対称性は

\[
[U,X]=0
\]

である。\(X\) の固有値 \(+1,-1\)
は異なるため、この可換条件を満たす作用素 \(U\)
は、対称部分空間と反対称部分空間を混合しない。

対称および反対称射影を

\[
P_s=
\frac12
\begin{pmatrix}
1&1\\
1&1
\end{pmatrix},
\qquad
P_a=
\frac12
\begin{pmatrix}
1&-1\\
-1&1
\end{pmatrix}
\]

とする。

このとき、

\[
P_s+P_a=I,
\qquad
P_sP_a=0
\]

であり、

\[
U=\lambda_sP_s+\lambda_aP_a
\]

と表せる。\(U\) がユニタリなら

\[
|\lambda_s|=|\lambda_a|=1
\]

である。全体位相 \(\lambda_s\) を除いて
\(\widetilde U:=\lambda_s^{-1}U\) とすれば、

\[
\boxed{
\widetilde U=P_s+\zeta P_a,
\qquad
\zeta:=\frac{\lambda_a}{\lambda_s},
\qquad
|\zeta|=1
}
\]

を得る。本稿の数値実装では初めから \(\lambda_s=1\) なので、以下では
\(\widetilde U\) を \(U_R\) と書く。このとき

\section{\$\$ U\_R}\label{-u_r}

\section{\textbackslash frac12 \textbackslash begin\{pmatrix\}
1+\textbackslash zeta\&1-\textbackslash zeta\textbackslash{}
1-\textbackslash zeta\&1+\textbackslash zeta
\textbackslash end\{pmatrix\}}\label{frac12-beginpmatrix-1zeta1-zeta-1-zeta1zeta-endpmatrix}

\textbackslash begin\{pmatrix\} r\&t\textbackslash{} t\&r
\textbackslash end\{pmatrix\}, \$\$

したがって

\[
\boxed{
r=\frac{1+\zeta}{2},
\qquad
t=\frac{1-\zeta}{2}
}
\]

である。すなわち交換振幅は、三角関数形を別途仮定しなくても、対称固有位相
\(1\) と反対称固有位相 \(\zeta\) の和・差として決まる。

さらに

\[
\boxed{
U_R^j=P_s+\zeta^jP_a
}
\]

と表せる。

この式は、モデルの複雑に見える時間発展が、実際には

\begin{itemize}
\tightlist
\item
  固定された対称成分、
\item
  単位円上を回転する反対称成分、
\end{itemize}

の二つから成ることを示す。

\subsection{4.3
観測量の周期性}\label{43-ux89b3ux6e2cux91cfux306eux5468ux671fux6027}

System Bの配分差は、一般に

\[
S_j=C\cos(j\omega+\varphi_0)
\]

と書ける。

ここで、

\[
\omega=\arg\lambda_a
\]

は作用素によって決まり、\(C\) および \(\varphi_0\)
は初期状態と観測関数によって決まる。

したがって、初期振幅、初期位相、内部倍音分布または波束形を変更しても、作用素固有の位相回転数
\(\omega(R)\)
は変わらない。変化するのは、その固有位相が選択した観測量にどの程度可視化されるかである。

\begin{center}\rule{0.5\linewidth}{0.5pt}\end{center}

\section{5. 有限位数再帰}\label{5-ux6709ux9650ux4f4dux6570ux518dux5e30}

\subsection{5.1
完全再帰条件}\label{51-ux5b8cux5168ux518dux5e30ux6761ux4ef6}

本稿の数値実装では対称固有値が厳密に1である。この位相固定された反復交換系が
\(n\) 回後に任意の初期状態を完全に復元する条件は、

\[
U_R^n=I
\]

である。

対称固有値は

\[
\lambda_s^n=1
\]

を満たすため、必要十分条件は反対称固有値について

\[
\lambda_a^n=1
\]

となることである。一般の交換対称 \(U\in U(2)\)
では、\(U^n=e^{i\gamma}I\)
は物理状態の全体位相を除く射影的再帰であり、\(\widetilde U=\lambda_s^{-1}U\)
に対する \(\widetilde U^n=I\) が以下の位相固定表現である。

\subsection{5.2
交換対称有限位数定理}\label{52-ux4ea4ux63dbux5bfeux79f0ux6709ux9650ux4f4dux6570ux5b9aux7406}

\textbf{定理5.1（交換対称有限位数写像の離散二乗重み）} 二チャネル作用素
\(U\in U(2)\) が次を満たすとする。

\begin{enumerate}
\def\labelenumi{\arabic{enumi}.}
\tightlist
\item
  チャネル交換作用素 \(X\) と可換である：\([U,X]=0\)。
\item
  全体位相を固定した作用素 \(\widetilde U:=\lambda_s^{-1}U\) を用いる。
\item
  \(\widetilde U\) は非自明な基本位数 \(n\ge3\)
  を持つ：\(\widetilde U^n=I\)。
\item
  反対称固有位相を \(\zeta=e^{-2\pi i m/n}\) とし、\(\gcd(m,n)=1\)
  とする。複素共役対 \(m\) と \(n-m\)
  は同じ重みを与えるため、代表元として \(1\le m<n/2\) を選ぶ。
\end{enumerate}

このとき

\section{\$\$ \textbackslash widetilde U}\label{-widetilde-u}

\section{P\_s+e\^{}\{-2\textbackslash pi i
m/n\}P\_a}\label{p_se-2pi-i-mnp_a}

\textbackslash begin\{pmatrix\} r\_\{n,m\}\&t\_\{n,m\}\textbackslash{}
t\_\{n,m\}\&r\_\{n,m\} \textbackslash end\{pmatrix\}, \$\$

\section{\$\$ r\_\{n,m\}}\label{-r_nm}

\section{\textbackslash frac\{1+e\^{}\{-2\textbackslash pi i
m/n\}\}\{2\}, \textbackslash qquad
t\_\{n,m\}}\label{frac1e-2pi-i-mn2-qquad-t_nm}

\textbackslash frac\{1-e\^{}\{-2\textbackslash pi i m/n\}\}\{2\}, \$\$

かつ、そのチャネル重みは

\section{\$\$ \textbackslash boxed\{
\textbar r\_\{n,m\}\textbar\^{}2}\label{-boxed-r_nm2}

\section{\textbackslash cos\^{}2\textbackslash left(\textbackslash frac\{\textbackslash pi
m\}\{n\}\textbackslash right), \textbackslash qquad
\textbar t\_\{n,m\}\textbar\^{}2}\label{cos2leftfracpi-mnright-qquad-t_nm2}

\textbackslash sin\^{}2\textbackslash left(\textbackslash frac\{\textbackslash pi
m\}\{n\}\textbackslash right) \} \$\$

となる。

\emph{証明.} \([U,X]=0\) より、\(U\) は \(X\)
の対称・反対称固有空間上で対角化される。全体位相を固定すると

\[
\widetilde U=P_s+\zeta P_a
\]

である。\(\widetilde U^n=I\) より \(\zeta^n=1\)。基本位数が \(n\) なので
\(\zeta=e^{-2\pi im/n}\)、\(\gcd(m,n)=1\)
と書ける。射影作用素をチャネル基底で展開すれば

\[
r_{n,m}=\frac{1+\zeta}{2},
\qquad
t_{n,m}=\frac{1-\zeta}{2}.
\]

したがって

\section{\$\$ \textbar r\_\{n,m\}\textbar\^{}2}\label{-r_nm2}

\section{\textbackslash frac\{(1+\textbackslash zeta)(1+\textbackslash zeta\^{}*)\}\{4\}}\label{frac1zeta1zeta4}

\section{\textbackslash frac\{1+\textbackslash cos(2\textbackslash pi
m/n)\}\{2\}}\label{frac1cos2pi-mn2}

\textbackslash cos\^{}2\textbackslash left(\textbackslash frac\{\textbackslash pi
m\}\{n\}\textbackslash right), \$\$

\section{\$\$ \textbar t\_\{n,m\}\textbar\^{}2}\label{-t_nm2}

\section{\textbackslash frac\{1-\textbackslash cos(2\textbackslash pi
m/n)\}\{2\}}\label{frac1-cos2pi-mn2}

\textbackslash sin\^{}2\textbackslash left(\textbackslash frac\{\textbackslash pi
m\}\{n\}\textbackslash right). \textbackslash qquad\textbackslash square
\$\$

\subsection{5.3
実装座標との一致と境界根}\label{53-ux5b9fux88c5ux5ea7ux6a19ux3068ux306eux4e00ux81f4ux3068ux5883ux754cux6839}

第3節の実装座標では

\[
\zeta=\lambda_a=-e^{2i\theta}
\]

であった。これを \(e^{-2\pi im/n}\) と同定すると

\section{\$\$ \textbackslash theta}\label{-theta}

\textbackslash frac\{\textbackslash pi\}\{2\}-\textbackslash frac\{\textbackslash pi
m\}\{n\} \textbackslash pmod\{\textbackslash pi\} \$\$

であり、

\section{\$\$
R=\textbar r\textbar\^{}2=\textbackslash sin\^{}2\textbackslash theta}\label{-rr2sin2theta}

\textbackslash cos\^{}2\textbackslash left(\textbackslash frac\{\textbackslash pi
m\}\{n\}\textbackslash right) \$\$

を再び得る。したがって三角関数形は定理の入力ではなく、射影分解から得られる係数を
\(R\) で表示した実装座標である。

定理5.1は非自明な内部根 \(0<R<1\) を対象とする。境界には

\[
\zeta=1:\quad U=I,\quad R=1,
\]

\[
\zeta=-1:\quad U=X,\quad R=0
\]

があり、それぞれ位数1および位数2の自明な閉包である。

\subsection{5.4
連続パラメータ内部の離散系列}\label{54-ux9023ux7d9aux30d1ux30e9ux30e1ux30fcux30bfux5185ux90e8ux306eux96e2ux6563ux7cfbux5217}

\(R\) は形式上、連続区間

\[
R\in[0,1]
\]

を動く。

しかし完全再帰する点は、整数対 \((n,m)\) によって

\section{\$\$ (n,m) \textbackslash longmapsto
R\_\{n,m\}}\label{-nm-longmapsto-r_nm}

\textbackslash cos\^{}2\textbackslash left(\textbackslash frac\{\textbackslash pi
m\}\{n\}\textbackslash right) \$\$

と分類される。

したがって、表面上連続な交換系の内部に、整数比位相によって番号付けされた可算な閉軌道系列が存在する。

有理数 \(m/n\)
は区間内に稠密であるため、高位数まで許せば有限位数根も稠密になる。一方、それぞれの根では

\[
U_R^n=I
\]

が厳密に成立し、近似的周期ではなく完全な離散閉軌道となる。

\begin{center}\rule{0.5\linewidth}{0.5pt}\end{center}

\section{6. 離散Born型重み}\label{6-ux96e2ux6563bornux578bux91cdux307f}

\subsection{6.1
標準二状態射影との形式的同型}\label{61-ux6a19ux6e96ux4e8cux72b6ux614bux5c04ux5f71ux3068ux306eux5f62ux5f0fux7684ux540cux578b}

標準的な二状態系では、

\section{\$\$
\textbar\textbackslash psi\textbackslash rangle}\label{-psirangle-1}

\textbackslash cos\textbackslash phi,\textbar A\textbackslash rangle +
e\^{}\{i\textbackslash chi\}\textbackslash sin\textbackslash phi,\textbar B\textbackslash rangle
\$\$

に対して、Born則は

\[
P(A)=\cos^2\phi,
\qquad
P(B)=\sin^2\phi
\]

を与える。

本研究の作用素をチャネル基底状態 \(|A\rangle=(1,0)^T\) に作用させると、

\section{\$\$
U\_\{n,m\}\textbar A\textbackslash rangle}\label{-u_nmarangle}

r\_\{n,m\}\textbar A\textbackslash rangle+t\_\{n,m\}\textbar B\textbackslash rangle
\$\$

となる。したがって、出力を同じA/B基底へ射影した二つの遷移重みは厳密に

\section{\$\$ W\_\{A\textbackslash to A\}}\label{-w_ato-a}

\section{\textbar\textbackslash langle
A\textbar U\_\{n,m\}\textbar A\textbackslash rangle\textbar\^{}2}\label{langle-au_nmarangle2}

\textbar r\_\{n,m\}\textbar\^{}2, \$\$

\section{\$\$ W\_\{A\textbackslash to B\}}\label{-w_ato-b}

\section{\textbar\textbackslash langle
B\textbar U\_\{n,m\}\textbar A\textbackslash rangle\textbar\^{}2}\label{langle-bu_nmarangle2}

\textbar t\_\{n,m\}\textbar\^{}2 \$\$

である。定理5.1より、

\section{\$\$ W\_\{A\textbackslash to A\}(n,m)}\label{-w_ato-anm}

\textbackslash cos\^{}2\textbackslash phi\_\{n,m\}, \$\$

\section{\$\$ W\_\{A\textbackslash to B\}(n,m)}\label{-w_ato-bnm}

\textbackslash sin\^{}2\textbackslash phi\_\{n,m\} \$\$

であり、

\[
\phi_{n,m}:=\frac{\pi m}{n}
\]

である。

したがって、基底入力に対する一回交換のチャネル射影重みは

\[
\boxed{
W_{A\to A}(n,m)=\cos^2\phi_{n,m},
\qquad
W_{A\to B}(n,m)=\sin^2\phi_{n,m}
}
\]

となり、標準二状態Born則と同型である。この等式は単なる数式の見た目の一致ではなく、同じ「複素振幅を観測基底へ射影し、その絶対値を二乗する」操作である。

ただし、一般の重ね合わせ入力 \(\alpha|A\rangle+\beta|B\rangle\)
に対するA出力は \(r\alpha+t\beta\)
であり、その重みには干渉項が含まれる。したがって \(R=|r|^2\)
は、任意入力状態に対するA観測確率そのものではなく、A基底入力からAチャネルへ残る遷移重みである。また、この射影重みを多数回試行の実験頻度として読むには、別途頻度則が必要である。

\subsection{6.2
何が入力され、何が出現したか}\label{62-ux4f55ux304cux5165ux529bux3055ux308cux4f55ux304cux51faux73feux3057ux305fux304b}

本結果を正確に評価するため、初期定義に含まれていた構造と、解析後に出現した構造を分ける。

\subsubsection{初期定義に含まれていたもの}\label{ux521dux671fux5b9aux7fa9ux306bux542bux307eux308cux3066ux3044ux305fux3082ux306e}

\begin{enumerate}
\def\labelenumi{\arabic{enumi}.}
\tightlist
\item
  二チャネル複素振幅。
\item
  ユニタリ散乱条件。
\item
  A/Bチャネル交換対称性。
\item
  A/B基底で振幅絶対値の二乗をチャネル重みとして読むこと。
\item
  数値実装上の座標としての振幅 \(r,t\) の三角関数表示。
\end{enumerate}

したがって、複素振幅の絶対値二乗が非負のチャネル重みになる構造そのものは、現行モデルに含まれていた。ただし定理5.1により、三角関数表示は独立な導出仮定ではなく、交換対称ユニタリ作用素の射影分解から再構成できる。

\subsubsection{初期定義に含まれていなかったもの}\label{ux521dux671fux5b9aux7fa9ux306bux542bux307eux308cux3066ux3044ux306aux304bux3063ux305fux3082ux306e}

\begin{enumerate}
\def\labelenumi{\arabic{enumi}.}
\tightlist
\item
  どの \(R\) が特別なピークを作るか。
\item
  どの位相角が完全再帰として選択されるか。
\item
  根が整数対 \((n,m)\) で分類されること。
\item
  選択された交換重みが \(\cos^2(\pi m/n)\) という離散系列を作ること。
\item
  局在性移乗モデルと灰色準安定性指標が、共通線形核の有限位数根と接続されること。
\end{enumerate}

したがって本研究の非自明な帰結は、二乗則そのものを無前提から生成したことではなく、

\[
\boxed{
\text{有限位数閉包}
\Rightarrow
\text{離散位相系列}
\Rightarrow
\text{離散Born型重み系列}
}
\]

である。

\subsection{6.3
確率公理より先に現れる閉包選択重み}\label{63-ux78baux7387ux516cux7406ux3088ux308aux5148ux306bux73feux308cux308bux9589ux5305ux9078ux629eux91cdux307f}

通常、\(\cos^2\phi\) は測定確率として導入される。

本研究では、同じ形が先に

\begin{quote}
二チャネル交換が有限回で完全に閉じるための反射・透過重み
\end{quote}

として現れる。

すなわち、

\[
\text{確率としての }\cos^2
\]

ではなく、まず

\[
\text{完全再帰可能な交換比としての }\cos^2
\]

が得られる。

この順序は、Born型重みを単なる外部確率公理ではなく、閉鎖位相系の再帰条件として読む可能性を示す。

ここでいう「選択」は動力学的吸引ではない。掃引中の \(R\)
は時間発展変数ではなく、各 \(R\)
を固定して軌道の閉包誤差を比較している。したがって有限位数根は

\[
R_j\longrightarrow R_{n,m}
\]

というアトラクターではなく、

\[
U_R^n=I
\]

を満たすパラメータとして閉包判定により選び出される。本稿ではこれを「閉包選択」または「再帰適合重み」と呼ぶ。強観測によるA/B動力学的選択は、System
BのD反作用という別機構である。

\subsection{6.4
離散系列と連続Born則}\label{64-ux96e2ux6563ux7cfbux5217ux3068ux9023ux7d9abornux5247}

本研究で直接得られたのは、

\[
\phi_{n,m}=\frac{\pi m}{n}
\]

に対応する離散系列である。

\(n\) を大きくすると有理位相は連続角度へ稠密に近づくため、

\[
\cos^2\left(\frac{\pi m}{n}\right)
\]

の集合も \([0,1]\) に稠密になる。

ただし、稠密であることと、連続Born則を物理的に導出したことは同義ではない。連続極限を物理法則として主張するには、

\begin{itemize}
\tightlist
\item
  どの位数までが実現可能か、
\item
  高位数根が雑音下でどの程度保持されるか、
\item
  観測時間およびコヒーレンス時間がどのように根を選別するか、
\item
  実験頻度が二乗重みに収束するか、
\end{itemize}

を別途示す必要がある。

\begin{center}\rule{0.5\linewidth}{0.5pt}\end{center}

\section{7.
波束収縮に見える挙動との関係}\label{7-ux6ce2ux675fux53ceux7e2eux306bux898bux3048ux308bux6319ux52d5ux3068ux306eux95a2ux4fc2}

\subsection{7.1
非ユニタリ収縮ではなく局在性再配分}\label{71-ux975eux30e6ux30cbux30bfux30eaux53ceux7e2eux3067ux306fux306aux304fux5c40ux5728ux6027ux518dux914dux5206}

System
Aにおいて、一方の波束が低局在状態から高局在状態へ変化する場合、部分観測では波束が収縮したように見える。

正規化前の交換核 \(U_R\)
はユニタリであり、線形更新だけを取り出せば全体系は

\section{\$\$
\textbar A\_j\textbar\^{}2+\textbar B\_j\textbar\^{}2}\label{-a_j2b_j2}

\textbackslash text\{constant\} \$\$

を保つ。一方、実装されたSystem Aの全更新は

\section{\$\$ F\_R(A,B)}\label{-f_rab}

\textbackslash left( \textbackslash mathcal N(rA+tB),
\textbackslash mathcal N(tA+rB) \textbackslash right) \$\$

であり、各チャネルを別々に正規化する。このため \(F_R\)
自体をユニタリ作用素と呼ぶことはできない。

したがって、線形交換段階で生じているのは、

\[
\text{全状態の消失を伴う収縮}
\]

ではなく、

\[
\text{保存された全状態内での局在性と内部構造の再配分}
\]

である。System
Aの数値結果として現れる局在性移乗は、この保存的再配分をチャネル別正規化後に読んだものである。有限位数根が全非線形写像
\(F_R\) の厳密周期根でもあるかは、\(U_R\)
の定理から自動的には従わず、System A固有の検証命題となる。

\subsection{7.2
観測断面による収縮様表示}\label{72-ux89b3ux6e2cux65adux9762ux306bux3088ux308bux53ceux7e2eux69d8ux8868ux793a}

観測者が全二チャネル状態ではなく一方のチャネルまたは特定の局在性指標だけを読む場合、交換核と正規化の合成軌道の一断面が、

\begin{itemize}
\tightlist
\item
  拡散状態から局在状態への変化、
\item
  中間状態からA/B状態への選択、
\item
  波束幅の急激な縮小、
\end{itemize}

として見える可能性がある。

したがって、本モデルは

\section{\$\$ \textbackslash boxed\{
\textbackslash text\{見かけの収縮\}}\label{-boxed-textux898bux304bux3051ux306eux53ceux7e2e}

\textbackslash text\{ユニタリ交換核\} +
\textbackslash text\{局在性再配分\} +
\textbackslash text\{チャネル別正規化\} +
\textbackslash text\{限定された観測断面\} \} \$\$

という作業仮説を与える。

これは標準量子測定の状態更新を再現したという主張ではない。しかし、非ユニタリな瞬間収縮を基本操作として置かなくても、収縮に似た観測結果を作り得る具体的な数理モデルである。

\subsection{7.3
有限位数根による再帰可能な局在交換}\label{73-ux6709ux9650ux4f4dux6570ux6839ux306bux3088ux308bux518dux5e30ux53efux80fdux306aux5c40ux5728ux4ea4ux63db}

交換重みが

\[
R=R_{n,m}
\]

であるとき、線形核 \(U_R\) は任意の入力を \(n\) 回後に厳密に復元する。

したがって正規化前の線形発展において途中で局在性が一方へ集中しても、その変化は不可逆吸収ではなく、閉軌道上の一局面である。チャネル別正規化を含むSystem
Aの全軌道については、この閉包がどの条件で保存されるかを数値的に区別する必要がある。

この構造は、

\begin{itemize}
\tightlist
\item
  局在化して見える時点、
\item
  再び非局在化する時点、
\item
  A/Bが反転する時点、
\item
  初期状態へ完全回帰する時点、
\end{itemize}

を同一の周期軌道として扱えることを意味する。

\begin{center}\rule{0.5\linewidth}{0.5pt}\end{center}

\section{8.
白猫・黒猫・灰色猫系との関係}\label{8-ux767dux732bux9ed2ux732bux7070ux8272ux732bux7cfbux3068ux306eux95a2ux4fc2}

\subsection{8.1
準安定二状態}\label{81-ux6e96ux5b89ux5b9aux4e8cux72b6ux614b}

System Bでは、配分差

\[
S_j=|A_j|^2-|B_j|^2
\]

がA/B二状態の優勢度を表す。

反対称固有位相が回転するため、一般に

\[
S_j=C\cos(j\omega+\varphi_0)
\]

となる。

したがって灰色状態は、静的にAとBが半分ずつ混合しただけの状態ではなく、A/B配分差が小振幅で周期的に交換される動的準安定状態として表現できる。

\subsection{8.2
弱読出しと強観測選択}\label{82-ux5f31ux8aadux51faux3057ux3068ux5f37ux89b3ux6e2cux9078ux629e}

System Bの一回更新は、無摂動交換核と観測写像を分けて

\[
\begin{aligned}
\Psi_{j+1}^{(C)}
&=\mathcal C\!\left(U_R\Psi_j^{(C)}\right),\\
\Psi_{j+1}^{(D)}
&=\mathcal D\!\left(U_R\Psi_j^{(D)}\right)
\end{aligned}
\]

と書ける。弱いC読出しでは、観測作用が小さく、反復軌道を大きく変えずに配分差を読む。強いD観測は、読出した
\(S_{D,j}\) を用いて

\section{\$\$ S\_\{j+1\}}\label{-s_j1-1}

S\_j+g\_D S\_\{D,j\}(1-S\_j\^{}2) \$\$

と更新する明示的な非線形反作用を持ち、A/Bいずれかの優勢状態へ押し出す。したがってA/B選択の方向と不可逆性は有限位数核
\(U_R\) だけからは生じず、\(\mathcal D\) の構造と結合強度 \(g_D\)
に依存する。

この構造は、弱測定および強測定の標準理論をそのまま実装したものではない。しかし、

\begin{itemize}
\tightlist
\item
  状態を保持しやすい読出し、
\item
  状態を変化させる選択操作、
\end{itemize}

を同一モデル内で区別する。

ここで、有限位数根の位置を決めるのは
\(U_R\)、灰色軌道をなるべく保持して読むのは
\(\mathcal C\)、A/Bへ選択するのは \(\mathcal D\)
である。この役割分担により、再帰と観測選択を同じ作用素の効果として混同せずに接続できる。

\subsection{8.3
Born型重みが現れた文脈}\label{83-bornux578bux91cdux307fux304cux73feux308cux305fux6587ux8108}

重要なのは、\(\cos^2\)
系列が抽象的な二状態ベクトルだけから発見されたのではないことである。

この系列は、無摂動交換核の有限位数を評価する数値実験の中で、

\begin{enumerate}
\def\labelenumi{\arabic{enumi}.}
\tightlist
\item
  A/B二状態、
\item
  その中間にある灰色準安定状態、
\item
  弱読出しによる保持、
\item
  別途実装された強観測反作用によるA/B選択、
\item
  長時間ドリフトの小さい閉軌道、
\end{enumerate}

を数値的に調べた系のピーク原因を解析した結果として得られた。

したがって、Born型二乗則が通常用いられる物理的文脈と類似した数理モデル群において、無摂動核の有限位数条件から同じ形の重みが現れ、その同じ核へC/D観測写像を接続できた点に、本研究の主要な意味がある。

\begin{center}\rule{0.5\linewidth}{0.5pt}\end{center}

\section{9.
数値ピークと解析根}\label{9-ux6570ux5024ux30d4ux30fcux30afux3068ux89e3ux6790ux6839}

\subsection{9.1
ピーク評価量}\label{91-ux30d4ux30fcux30afux8a55ux4fa1ux91cf}

System Bでは、長さ \(k\) の時間列に対して、

\section{\$\$ \textbackslash overline S\_k}\label{-overline-s_k}

\textbackslash frac1k\textbackslash sum\_\{j=0\}\^{}\{k-1\}S\_j, \$\$

\section{\$\$ A\_k}\label{-a_k}

\textbackslash frac\{\textbackslash max\_\{j\textless k\}S\_j-\textbackslash min\_\{j\textless k\}S\_j\}\{2\},
\$\$

\section{\$\$ D\_k}\label{-d_k}

\subsection{\textbackslash left\textbar{} \textbackslash overline
S\_\{k,\textbackslash mathrm\{second\}\}}\label{left-overline-s_kmathrmsecond}

\textbackslash overline S\_\{k,\textbackslash mathrm\{first\}\}
\textbackslash right\textbar{} \$\$

を評価した。

有限位数根の直接確認に用いた主条件は、初期相対位相

\[
\varphi_0=0\ \text{または}\ \pi,
\]

コード上の初期差 \(s_0=0.01\)、したがって配分差

\[
S_0=|a_0|^2-|b_0|^2=0.02,
\]

目標振幅 \(C=0.02\) である。この条件に対し、gray errorを

\section{\$\$ \textbackslash varepsilon\_k}\label{-varepsilon_k}

\textbar\textbackslash overline S\_k\textbar{} +
\textbar A\_k-C\textbar{} + D\_k \$\$

と定義し、深度を

\[
d_k=-\log_{10}\varepsilon_k
\]

とした。

ここで評価しているのは、Born確率との一致ではない。

\begin{itemize}
\tightlist
\item
  平均が中立であること、
\item
  振動幅が保存されること、
\item
  前半と後半でドリフトしないこと、
\end{itemize}

を評価している。

\subsection{9.2 偶数位数根}\label{92-ux5076ux6570ux4f4dux6570ux6839}

主条件 \(\varphi_0=0\) または \(\pi\) では、定理5.1の根における配分差は

\section{\$\$ S\_j}\label{-s_j}

C\textbackslash cos\textbackslash left(\textbackslash frac\{2\textbackslash pi
m\}\{n\}j\textbackslash right), \textbackslash qquad C=0.02 \$\$

となる。基本位数 \(n\) が偶数で \(\gcd(m,n)=1\) なら \(m\)
は奇数であるため、

\[
S_0=C,
\qquad
S_{n/2}=C\cos(\pi m)=-C
\]

が格子点上で厳密に現れる。また一周期の円分位相和は0である。

したがって \(j=0,\ldots,2n-1\) の二完全周期、すなわち \(2n\)
個の標本を取れば、現在のgray metricについて

\[
\overline S=0,
\qquad
A=C,
\qquad
D=0
\]

が厳密に成立する。これは「偶数位数なら任意の初期状態と任意の評価窓でgray
errorが0になる」という主張ではない。作用素の再帰 \(U_R^n=I\)
は初期状態によらないが、上の振幅条件と誤差0は、\(\varphi_0=0\) または
\(\pi\)、\(S_0=C=0.02\)、二完全周期標本という実験条件に対する厳密結果である。

したがって、

\[
\varepsilon=0
\]

となり、理想演算では

\[
d\rightarrow+\infty
\]

である。

発散するのは物理的エネルギーまたは振幅ではなく、完全周期軌道に対する誤差指標の逆対数である。有限精度計算では丸め誤差によって有限の深度として表示される。

\subsection{9.3
奇数位数根と観測量依存性}\label{93-ux5947ux6570ux4f4dux6570ux6839ux3068ux89b3ux6e2cux91cfux4f9dux5b58ux6027}

基本位数 \(n\) が奇数の場合でも、

\[
U_R^n=I
\]

は成立する。

しかし有限標本列が正負両方の極値を格子点として含まない場合、現在の振幅評価には小さな欠損が残る。

したがって、

\[
\boxed{
\text{作用素の完全再帰}
\neq
\text{特定の観測量が示す無限深ピーク}
}
\]

である。

これは、Born型重みの候補系列と、その系列が特定の実験でどの程度可視化されるかを区別する必要を示す。

\subsection{9.4
観測された二つの主ピーク}\label{94-ux89b3ux6e2cux3055ux308cux305fux4e8cux3064ux306eux4e3bux30d4ux30fcux30af}

全域掃引と局所精密掃引で追跡した二つの主ピークは、定理5.1の次の根に一致する{[}4{]}。

{\def\LTcaptype{none} % do not increment counter
\begin{longtable}[]{@{}lrrrr@{}}
\toprule\noalign{}
根 & 基本位数 \(n\) & \(m\) & \(R_{n,m}=\lvert r\rvert^2\) &
\(T_{n,m}=\lvert t\rvert^2\) \\
\midrule\noalign{}
\endhead
\bottomrule\noalign{}
\endlastfoot
\(R_{124,23}\) & 124 & 23 & 0.697177927556659 & 0.302822072443341 \\
\(R_{122,23}\) & 122 & 23 & 0.688363946817593 & 0.311636053182407 \\
\end{longtable}
}

両者は偶数基本位数かつ \(m=23\)
が奇数である。このため、§9.2の主条件では正負両極値が標本格子上に現れ、理想演算でgray
errorが0となる。数値掃引で見えた鋭いピークは、この偶数位数円分再帰を有限精度で読んだものである。

ここで同定されたのは、微細構造定数そのものではなく、交換作用素の有限位数根である。これらが物理定数近傍に位置する理由は本稿の定理からは導かれず、先行報告{[}4{]}に残された別問題である。

\begin{center}\rule{0.5\linewidth}{0.5pt}\end{center}

\section{10.
本研究が示すこと}\label{10-ux672cux7814ux7a76ux304cux793aux3059ux3053ux3068}

\subsection{10.1
導出済みの数理的帰結}\label{101-ux5c0eux51faux6e08ux307fux306eux6570ux7406ux7684ux5e30ux7d50}

本研究から厳密に導けるのは次である。

\begin{enumerate}
\def\labelenumi{\arabic{enumi}.}
\tightlist
\item
  交換対称な二チャネルユニタリ作用素は、全体位相を除いて
\end{enumerate}

\[
U=P_s+\zeta P_a,
\qquad |\zeta|=1
\]

と一意に表される。 2. チャネル基底における交換振幅は

\[
r=\frac{1+\zeta}{2},
\qquad
t=\frac{1-\zeta}{2}
\]

であり、三角関数形を独立に仮定する必要はない。 3.
位相固定後の完全再帰条件 \(U^n=I\)
は、反対称固有位相が根1になることと同値である。 4. 非自明な有限位数根
\(\zeta=e^{-2\pi im/n}\) におけるA基底入力のA/B遷移重みは

\[
W_{A\to A}(n,m)=\cos^2\left(\frac{\pi m}{n}\right)
\]

で与えられる。 5. 対応する相補重みは

\[
W_{A\to B}(n,m)=\sin^2\left(\frac{\pi m}{n}\right)
\]

である。 6.
これらは、A/B基底への振幅射影の絶対値二乗として、標準二状態Born則と同じ数学的操作および
\(\cos^2/\sin^2\) 形を持つ。 7.
この離散系列は、Born則を評価関数として入力した結果ではなく、交換対称性、ユニタリ性および有限位数再帰から得られた。
8. 上の定理は共通線形核 \(U\) に対する結果であり、System
Aのチャネル別正規化およびSystem
BのC/D反作用を含む全写像の同一性を主張しない。

\subsection{10.2
物理的に示唆されること}\label{102-ux7269ux7406ux7684ux306bux793aux5506ux3055ux308cux308bux3053ux3068}

本研究は次の可能性を示唆する。

\begin{quote}
Born型二乗重みは、交換対称な閉鎖二チャネル波動系における有限回帰可能な固有位相を、チャネル基底へ射影した遷移重みとして現れる。
\end{quote}

さらに、

\begin{quote}
波束収縮に見える局在性集中は、ユニタリ交換核、局在性再配分、チャネル別正規化および限定された観測断面の合成として表現できる。強観測による二状態選択は、同じ交換核へ非線形D反作用を接続した別の動力学として表現できる。
\end{quote}

\subsection{10.3
本研究がまだ示していないこと}\label{103-ux672cux7814ux7a76ux304cux307eux3060ux793aux3057ux3066ux3044ux306aux3044ux3053ux3068}

本稿は次を証明していない。

\begin{enumerate}
\def\labelenumi{\arabic{enumi}.}
\tightlist
\item
  標準量子力学のBorn則全体の導出。
\item
  単一試行における確率的結果の発生機構。
\item
  多数回試行の頻度が \(\cos^2\phi\) に収束すること。
\item
  任意次元ヒルベルト空間における一般射影測度の一意性。
\item
  実際の量子測定における不可逆性。
\item
  観測者、環境またはデコヒーレンスの完全な動力学。
\item
  System Aの局在性移乗が標準的な波動関数収縮と同一であること。
\item
  チャネル別正規化を含むSystem
  Aの全写像が、線形核と同じ有限位数で厳密に閉じること。
\item
  一般の \(U(2)\)
  作用素で、交換対称性なしに同じチャネル重みが一意に決まること。
\end{enumerate}

したがって、本稿では「Born則を導出した」ではなく、

\begin{quote}
有限位数回帰から離散Born型重みが出現した
\end{quote}

と表現する。

\begin{center}\rule{0.5\linewidth}{0.5pt}\end{center}

\section{11.
標準Born則との位置関係}\label{11-ux6a19ux6e96bornux5247ux3068ux306eux4f4dux7f6eux95a2ux4fc2}

\subsection{11.1
Bornの統計解釈}\label{111-bornux306eux7d71ux8a08ux89e3ux91c8}

Bornは1926年、量子散乱において波動関数の振幅から事象の確率を読む統計解釈を導入した{[}5{]}。現代的には、規格化状態
\(|\psi\rangle\) と射影 \(P_A\) に対して、

\[
P(A)=\langle\psi|P_A|\psi\rangle
\]

と表される。

純粋状態かつ一次元射影では、

\[
P(A)=|\langle A|\psi\rangle|^2
\]

である。

\subsection{11.2
Gleasonの定理との違い}\label{112-gleasonux306eux5b9aux7406ux3068ux306eux9055ux3044}

Gleasonの定理は、次元3以上のヒルベルト空間における直交射影への加法的確率測度が、密度作用素による

\[
\mu(P)=\operatorname{Tr}(\rho P)
\]

の形になることを示す{[}6{]}。

これはBorn型測度の数学的一意性に関する強い結果である。

通常のGleason定理は二次元ヒルベルト空間を直接は含まない。一方、Buschは射影測定より広い効果作用素への加法性を用い、密度作用素表示を得る一般化Gleason型結果を示している{[}7{]}。したがって、二状態系にもBorn型測度を特徴づける既知の数学的枠組みは存在する。

一方、本研究はその一般性を扱わない。本稿の対象は二チャネル反復交換系であり、問題は

\begin{quote}
なぜ特定の二乗射影値が有限回帰可能な交換重みとして選ばれるのか
\end{quote}

である。

したがって両者の役割は異なる。

\begin{itemize}
\tightlist
\item
  Gleason型結果：整合的な確率測度の形を制約する。
\item
  本研究：反復交換力学の閉包条件が選ぶ離散二乗重みを示す。
\end{itemize}

\subsection{11.3
他のBorn則導出との違い}\label{113-ux4ed6ux306ebornux5247ux5c0eux51faux3068ux306eux9055ux3044}

Zurekのenvarianceによる議論は、複合系のエンタングルメント対称性からSchmidt成分の確率
\(p_k\propto|\psi_k|^2\)
を導くことを目的とする{[}8{]}。これは「なぜ振幅二乗を確率として読むか」という問題へ直接取り組む先行研究である。

本研究の問いは異なる。ここでは確率概念を出発点にせず、交換対称な一体の二チャネル作用素に対して

\[
U^n=I
\]

となる固有位相を分類し、そのチャネル射影重みが \(\cos^2/\sin^2\)
になることを示す。したがって本稿はenvarianceやGleason型測度定理の代替ではなく、有限位数回帰からBorn型射影重みが生じる別の数理経路を提示する。

\subsection{11.4 exact
recurrence研究との接点}\label{114-exact-recurrenceux7814ux7a76ux3068ux306eux63a5ux70b9}

Anandらは、有限次元Floquet系のexactかつ初期状態に依存しない再帰を、ユニタリスペクトルの円分構造から分類している{[}9{]}。本稿の

\[
\zeta^n=1
\]

という条件は、このexact
recurrenceの最小二チャネル実現に当たる。先行研究が一般のFloquetスペクトルについて再帰時間の算術的分類を行うのに対し、本稿は交換対称性によって固有ベクトルを対称・反対称モードへ固定し、その円分固有位相をA/Bチャネルへ戻すと

\section{\$\$ \textbackslash left\textbar{}
\textbackslash frac\{1\textbackslash pm\textbackslash zeta\}\{2\}
\textbackslash right\textbar\^{}2}\label{-left-frac1pmzeta2-right2}

\textbackslash cos\^{}2\textbackslash phi,\textbackslash{}
\textbackslash sin\^{}2\textbackslash phi \$\$

が現れることを示す。この「円分再帰からチャネル二乗重みへの写像」が、本稿固有の焦点である。

\subsection{11.5
二状態系に限定した意味}\label{115-ux4e8cux72b6ux614bux7cfbux306bux9650ux5b9aux3057ux305fux610fux5473}

本研究は二状態系であるため、Gleasonの定理を直接適用してBorn則を導出するものではない。

しかし二状態系は、

\begin{itemize}
\tightlist
\item
  干渉、
\item
  スピン二準位、
\item
  二経路測定、
\item
  二モード光学、
\item
  A/B観測選択、
\end{itemize}

の最小模型である。

その最小模型において、再帰可能な交換重みが

\[
\cos^2\phi,
\qquad
\sin^2\phi
\]

として現れることは、一般理論へ進む前の成立範囲が明確な結果である。

\begin{center}\rule{0.5\linewidth}{0.5pt}\end{center}

\section{12.
反証可能な検証課題}\label{12-ux53cdux8a3cux53efux80fdux306aux691cux8a3cux8ab2ux984c}

\subsection{\texorpdfstring{12.1 交換対称部分群と一般
\(U(2)\)}{12.1 交換対称部分群と一般 U(2)}}\label{121-ux4ea4ux63dbux5bfeux79f0ux90e8ux5206ux7fa4ux3068ux4e00ux822c-u2}

定理5.1により、\(\cos^2\)
形が第3節の散乱係数表示だけに由来する可能性は排除された。必要なのは特定の三角関数座標ではなく、

\[
[U,X]=0
\]

というA/B交換対称性である。

一方、一般の二準位ユニタリ作用素は、全体位相を含めて

\section{\$\$ U}\label{-u}

e\^{}\{i\textbackslash alpha\} \textbackslash left(
\textbackslash cos\textbackslash beta,I
-i\textbackslash sin\textbackslash beta,\textbackslash boldsymbol
n\textbackslash cdot\textbackslash boldsymbol\textbackslash sigma
\textbackslash right), \textbackslash qquad
\textbar\textbackslash boldsymbol n\textbar=1 \$\$

と書ける。有限位数条件は固有位相 \(\beta\)
を有理角へ量子化するが、A基底からB基底への遷移重みは

\section{\$\$ \textbar\textbackslash langle
B\textbar U\textbar A\textbackslash rangle\textbar\^{}2}\label{-langle-buarangle2}

\textbackslash sin\^{}2\textbackslash beta,(n\_x\^{}2+n\_y\^{}2) \$\$

であり、固有ベクトル軸 \(\boldsymbol n\)
が残る。したがって、有限位数条件だけでは一般 \(U(2)\) のA/B重みを一意に
\(\sin^2\beta\) へ固定できない。

本稿の定理が適用されるのは、交換作用素 \(X=\sigma_x\)
の中心化部分群、すなわち \(\boldsymbol n\)
が交換軸へ固定され、対称・反対称固有ベクトルがA/B基底に対して等重みとなる場合である。これは定理の欠陥ではなく、離散Born型重みを生む正確な対称性条件の同定である。

反証可能な次の検証は、\([U,X]\ne0\) となる制御摂動を入れ、上の軸因子
\(n_x^2+n_y^2\)
に従ってピーク重みが変化するかを測ることである。変化しなければ、本稿の射影機構だけでは説明できない追加構造が存在する。

\subsection{12.2
頻度則の検証}\label{122-ux983bux5ea6ux5247ux306eux691cux8a3c}

現行結果は閉軌道重みを与えるが、測定頻度を直接生成していない。

次段階では、同一のA基底入力 \(|A\rangle\)
を用意し、有限位数核と明示した測定・選択写像を通す多数試行について、A/B出力頻度

\[
f_{A\to A},
\qquad
f_{A\to B}
\]

を測定し、

\[
f_{A\to A}\rightarrow W_{A\to A}(n,m)=|r_{n,m}|^2,
\qquad
f_{A\to B}\rightarrow W_{A\to B}(n,m)=|t_{n,m}|^2
\]

となるかを調べる必要がある。

これが成立すれば、有限位数交換重みと実験頻度としてのBorn則の間に、より直接的な接続が得られる。一般の重ね合わせ入力については、\(r\alpha+t\beta\)
の干渉項を含む頻度則を別に検証する。

\subsection{12.3
観測摂動と不可逆性}\label{123-ux89b3ux6e2cux6442ux52d5ux3068ux4e0dux53efux9006ux6027}

完全ユニタリ系は再帰可能であり、原理的に不可逆ではない。

現実の観測選択へ接続するには、

\begin{itemize}
\tightlist
\item
  観測装置自由度、
\item
  環境自由度、
\item
  位相雑音、
\item
  チャネル損失、
\item
  読出し後のフィードバック、
\end{itemize}

を導入し、再帰可能な軌道が実効的にA/Bどちらかへ固定される条件を調べる必要がある。

\subsection{12.4
ノイズ下の有限位数選択}\label{124-ux30ceux30a4ux30baux4e0bux306eux6709ux9650ux4f4dux6570ux9078ux629e}

根からのずれを

\[
R=R_{n,m}+\delta R
\]

とする。

\(n\) 回後の作用素残差

\[
\mathcal E_n(R)
:=
\|U_R^n-I\|
\]

を測定すれば、各有限位数根の共鳴幅を定義できる。

さらに反復ごとの位相雑音を

\[
\theta_j=\theta+\xi_j
\]

として、

\[
U_{j+n-1}\cdots U_j
\]

の再帰忠実度を評価する。

これにより、数学的には稠密に存在する根のうち、どの低位数根が物理的に頑健かを選別できる。

\subsection{12.5 System
Aの正規化監査}\label{125-system-aux306eux6b63ux898fux5316ux76e3ux67fb}

System Aの実装にはチャネル別正規化が含まれており、全更新写像は線形な
\(U_R\otimes I\) ではなく、状態依存の非線形写像である。

したがって、

\begin{enumerate}
\def\labelenumi{\arabic{enumi}.}
\tightlist
\item
  正規化なし、
\item
  二チャネル全体を一括正規化、
\item
  各チャネルを個別正規化、
\end{enumerate}

の三条件を比較し、有限位数根の位置、局在性移乗およびBorn型重みの可視性を分離する必要がある。

この監査は、共通線形核について導出された有限位数構造と、System
Aの非線形な波束局在性挙動を厳密に接続するために必要である。現行実装がチャネル別正規化を持つこと自体は確認済みであり、未確認なのは、その正規化が線形核の有限位数根をどの条件で保存・移動・消失させるかである。

\begin{center}\rule{0.5\linewidth}{0.5pt}\end{center}

\section{13. 考察}\label{13-ux8003ux5bdf}

\subsection{13.1
本研究の主要な気づき}\label{131-ux672cux7814ux7a76ux306eux4e3bux8981ux306aux6c17ux3065ux304d}

本研究の主要な気づきは、連続な交換係数 \(R\)
を持つ一見連続的な波動系の内部に、整数対 \((n,m)\)
で分類される完全な離散閉軌道が存在し、その重みが交換対称性から直接決まることである。

しかも、その離散閉軌道の交換重みは

\[
R_{n,m}=\cos^2\left(\frac{\pi m}{n}\right)
\]

となる。

交換対称性とユニタリ性は

\[
U=P_s+\zeta P_a
\]

を与え、有限位数閉包は \(\zeta=e^{-2\pi im/n}\) を与える。したがって、

\[
\text{離散位相閉包}
\longrightarrow
\text{振幅射影}
\longrightarrow
\text{二乗重み}
\]

という構造が成立する。\(\cos^2\)
は仮定した散乱座標の名残ではなく、二つの射影固有位相の和・差をA/B基底へ戻したときの干渉二乗である。

\subsection{13.2
Born型重みの有限位数起源}\label{132-bornux578bux91cdux307fux306eux6709ux9650ux4f4dux6570ux8d77ux6e90}

Born則では二乗重みは観測確率として与えられる。

本研究では同じ二乗重みが、有限回で自己同一状態へ戻るための交換比として現れる。

したがって、本モデル内では次の命題が定理として得られる。

\begin{quote}
交換対称な閉鎖二チャネル系では、再帰可能な円分固有位相をA/Bチャネルへ射影した遷移重みが、離散Born型二乗系列となる。
\end{quote}

この導出では、確率は最初から無構造な乱数として置かれず、

\begin{enumerate}
\def\labelenumi{\arabic{enumi}.}
\tightlist
\item
  閉じた位相関係、
\item
  有限回帰可能な整数比、
\item
  観測基底への射影、
\item
  射影成分の二乗、
\end{enumerate}

という順序で現れる。未解決なのは、この遷移重みが物理的な多数回測定頻度を支配する機構であり、離散二乗重みそのものの導出ではない。

\subsection{13.3
収縮・準安定・選択の共通核}\label{133-ux53ceux7e2eux6e96ux5b89ux5b9aux9078ux629eux306eux5171ux901aux6838}

System AおよびSystem Bは、異なる現象を調べるために構築された。

しかし解析後には、

\begin{itemize}
\tightlist
\item
  波束収縮に見える局在性集中、
\item
  A/B二状態の周期交換、
\item
  灰色準安定状態、
\item
  弱読出しによる保持、
\item
  強観測による選択、
\item
  離散Born型重み、
\end{itemize}

が、同じ対称・反対称固有モード分解を持つ無摂動核へ接続された。ただし全更新写像は同一ではない。System
Aの局在性表示にはチャネル別正規化 \(\mathcal N_A\) が、System
Bの弱読出しには \(\mathcal C\) が、強観測選択には非線形反作用
\(\mathcal D\) が加わる。

このことは、量子測定において別々の問題として語られる現象を、最小二チャネル模型では「共通の交換位相核」と「異なる後段写像」の組合せとして比較できることを示す。根の位置、局在性表示、弱読出し、強選択を一つの効果へ混同せず、同じ計算模型内で接続できる点が重要である。

\subsection{13.4
厳密な中心結論}\label{134-ux53b3ux5bc6ux306aux4e2dux5fc3ux7d50ux8ad6}

本稿の厳密な中心結論は次である。

\begin{quote}
波束局在性移乗、準安定二状態、弱読出しおよび強観測選択を調べるために独立に構成されたモデル群の共通線形核は、交換対称な二チャネルユニタリ作用素である。この核の有限位数閉包は、Born則を探索目標として与えることなく、\(\cos^2(\pi m/n)\)
と \(\sin^2(\pi m/n)\) という離散チャネル射影重みを必然的に与える。
\end{quote}

これはBorn則全体の導出ではない。しかし、交換対称性、ユニタリ性、有限位数閉包およびチャネル射影から離散Born型二乗重みが生じることを、再現可能なモデルと解析式で示す。頻度則、System
Aの非線形閉包、およびD観測と標準的測定動力学との対応は、この定理とは分離された次の課題である。

\begin{center}\rule{0.5\linewidth}{0.5pt}\end{center}

\section{14. 結論}\label{14-ux7d50ux8ad6}

本研究は、チャネル交換作用素 \(X\)
と可換な二チャネルユニタリ作用素を解析した。交換対称性により、全体位相を固定した作用素は

\[
U=P_s+\zeta P_a,
\qquad |\zeta|=1
\]

と表される。チャネル基底へ戻すと、振幅は

\[
r=\frac{1+\zeta}{2},
\qquad
t=\frac{1-\zeta}{2}
\]

である。有限位数条件 \(U^n=I\) は反対称固有位相を

\[
\zeta=e^{-2\pi im/n}
\]

へ量子化し、有限回で完全再帰するチャネル射影重み

\[
\boxed{
|r_{n,m}|^2=\cos^2\left(\frac{\pi m}{n}\right)
}
\]

および相補重み

\[
\boxed{
|t_{n,m}|^2=\sin^2\left(\frac{\pi m}{n}\right)
}
\]

を与える。この \(\cos^2/\sin^2\)
形式は、A基底入力をA/B基底へ射影した振幅の絶対値二乗として、標準量子力学における二状態Born則と同じ数学的構造である。三角関数形は独立な仮定ではなく、交換対称射影
\(P_s,P_a\) と円分固有位相の干渉から導かれる。

本研究の重要点は、この形を再現するためにモデルを設計したのではないことである。出発点は、波束局在性の移乗、収縮に見える局在化、白猫・黒猫・灰色猫準安定界面、弱い読出しおよび強観測選択を調べる二つの数値モデルであった。観測された鋭いピークの原因を後から固有値解析した結果、有限位数根と離散Born型重みが得られ、さらに本稿で三角関数パラメータに依存しない交換対称有限位数定理へ高められた。

同時に、実装の役割も分離された。根の位置を与えるのは共通線形核
\(U\)、System Aの局在性表示を作るのはチャネル別正規化を含む
\(\mathcal N_A\circ U\)、System Bの弱読出しと強選択を作るのはそれぞれ
\(\mathcal C\circ U\) と \(\mathcal D\circ U\) である。

したがって、本稿の結論はBorn則の完全導出ではない。

本稿が示したのは、

\[
\boxed{
\begin{gathered}
\text{交換対称ユニタリ核}
+\text{円分有限位数閉包}
+\text{A/Bチャネル射影}\\
\Rightarrow\text{離散Born型二乗重み}
\end{gathered}
}
\]

という成立条件の明確な数理的接続である。

次段階では、交換対称性を破る一般 \(U(2)\)
制御実験、測定頻度の再現、観測装置および環境自由度の導入、ノイズ下での根の選別、ならびにSystem
Aの正規化条件別監査を行う。

\begin{center}\rule{0.5\linewidth}{0.5pt}\end{center}

\section{参考文献}\label{ux53c2ux8003ux6587ux732e}

\subsection{自己引用}\label{ux81eaux5df1ux5f15ux7528}

\begin{enumerate}
\def\labelenumi{\arabic{enumi}.}
\tightlist
\item
  木原範昭,
  「交換干渉散乱行列フェルミオン的衝突における低局在性・倍音移乗読出し実験仕様書
  v1」, 2026.
\item
  木原範昭,
  「交換干渉散乱行列フェルミオン的衝突における加速度基底と局在性交換予備実験総括
  v1」, Zenodo Concept DOI: 10.5281/zenodo.21333766, 2026.
\item
  木原範昭,
  「白猫黒猫灰色猫準安定界面におけるC弱読出しとD強観測選択予備実験総括
  v1」, Zenodo Concept DOI: 10.5281/zenodo.21353208, 2026.
\item
  木原範昭,
  「反復交換散乱における有限位数共鳴の発見――微細構造定数137・128近傍ピークの原因特定と再現可能な波束数理モデル」,
  Version DOI: 10.5281/zenodo.21421367, Concept DOI:
  10.5281/zenodo.21421366, 2026.
\end{enumerate}

\subsection{外部引用}\label{ux5916ux90e8ux5f15ux7528}

\begin{enumerate}
\def\labelenumi{\arabic{enumi}.}
\setcounter{enumi}{4}
\tightlist
\item
  Max Born, ``Zur Quantenmechanik der Stoßvorgänge,'' \emph{Zeitschrift
  für Physik} \textbf{37}, 863--867 (1926). DOI: 10.1007/BF01397477.
\item
  Andrew M. Gleason, ``Measures on the Closed Subspaces of a Hilbert
  Space,'' \emph{Journal of Mathematics and Mechanics} \textbf{6},
  885--893 (1957). DOI: 10.1512/iumj.1957.6.56050.
\item
  Paul Busch, ``Quantum States and Generalized Observables: A Simple
  Proof of Gleason\textquotesingle s Theorem,'' \emph{Physical Review
  Letters} \textbf{91}, 120403 (2003). DOI:
  10.1103/PhysRevLett.91.120403.
\item
  Wojciech H. Zurek, ``Probabilities from Entanglement:
  Born\textquotesingle s Rule \(p_k=|\psi_k|^2\) from Envariance,''
  \emph{Physical Review A} \textbf{71}, 052105 (2005). DOI:
  10.1103/PhysRevA.71.052105.
\item
  Amit Anand, Dinesh Valluri, Jack Davis, and Shohini Ghose, ``Quantum
  Recurrences and the Arithmetic of Floquet Dynamics,'' \emph{Quantum}
  \textbf{10}, 2074 (2026). DOI: 10.22331/q-2026-04-20-2074.
\end{enumerate}

\begin{center}\rule{0.5\linewidth}{0.5pt}\end{center}

\section{付録A.
主張の分類}\label{ux4ed8ux9332a-ux4e3bux5f35ux306eux5206ux985e}

{\def\LTcaptype{none} % do not increment counter
\small
\begin{longtable}[]{@{}p{0.54\linewidth}p{0.39\linewidth}@{}}
\toprule\noalign{}
主張 & 分類 \\
\midrule\noalign{}
\endhead
\bottomrule\noalign{}
\endlastfoot
交換対称ユニタリ作用素の位相固定表示 \(U=P_s+\zeta P_a\) &
定理5.1の導出済み帰結 \\
有限位数条件 \(U^n=I\Leftrightarrow\zeta^n=1\) &
定理5.1の導出済み帰結 \\
\(r=(1+\zeta)/2\), \(t=(1-\zeta)/2\) & 射影分解からの導出済み帰結 \\
\(\lvert r_{n,m}\rvert^2=\cos^2(\pi m/n)\) & 導出済み帰結 \\
\(\lvert t_{n,m}\rvert^2=\sin^2(\pi m/n)\) & 導出済み帰結 \\
三角関数形は独立仮定ではなく実装座標として回収される & 導出済み帰結 \\
基底入力に対する有限位数重みと二状態Born射影の同型 & 導出済み比較 \\
\(\cos^2\) 系列を探索目標として入力していない &
実験・開発履歴上の事実 \\
System AとSystem Bが同じ正規化前・無摂動線形核を共有する &
ソースコードで確認済みの実装事実 \\
System Aがチャネル別正規化を持ち、全写像が非線形である &
ソースコードで確認済みの実装事実 \\
System Aの全非線形写像が線形核と同じ位数で厳密再帰する & 未検証 \\
System BのD選択が明示的な非線形反作用を持つ &
ソースコードで確認済みの実装事実 \\
偶数位数根でgray errorが0となる & \(\varphi_0=0\) または
\(\pi\)、\(S_0=C=0.02\)、\(2n\) 標本に限定した導出済み帰結 \\
無限深ピーク & 誤差指標 \(-\log_{10}\varepsilon\)
の数学的発散。エネルギー発散ではない \\
一般 \(U(2)\) では有限位数だけで同じA/B重みは決まらない &
§12.1の導出済み帰結 \\
波束収縮に見える挙動が交換核・正規化・部分観測で表現できる & 現行System
Aモデル内の数理的知見 \\
標準量子測定の波束収縮を説明した & 未導出 \\
Born則全体を導出した & 未導出 \\
多数回測定頻度が \(\cos^2\) に収束する & 未検証 \\
離散Born型チャネル重みが有限位数位相閉包に由来する &
交換対称二チャネル核について導出済み \\
物理的Born頻度則の起源が有限位数位相閉包である & 今後検証する物理仮説 \\
\end{longtable}
}

\begin{center}\rule{0.5\linewidth}{0.5pt}\end{center}

\section{付録B.
最小再現式}\label{ux4ed8ux9332b-ux6700ux5c0fux518dux73feux5f0f}

本稿の中心結果は、三角関数表示を仮定せず、次の式だけから再現できる。

\subsection{B.1
交換対称射影}\label{b1-ux4ea4ux63dbux5bfeux79f0ux5c04ux5f71}

\[
X=
\begin{pmatrix}
0&1\\
1&0
\end{pmatrix},
\]

\[
P_s=\frac{I+X}{2},
\qquad
P_a=\frac{I-X}{2}.
\]

\subsection{B.2
交換対称ユニタリ作用素}\label{b2-ux4ea4ux63dbux5bfeux79f0ux30e6ux30cbux30bfux30eaux4f5cux7528ux7d20}

\[
[U,X]=0,
\qquad
U^\dagger U=I
\]

ならば、全体位相を固定して

\[
\boxed{
U=P_s+\zeta P_a,
\qquad |\zeta|=1
}
\]

と書ける。

\subsection{B.3
有限位数条件}\label{b3-ux6709ux9650ux4f4dux6570ux6761ux4ef6}

非自明な基本位数 \(n\ge3\) に対し、

\[
U^n=I
\quad\Longleftrightarrow\quad
\zeta^n=1.
\]

したがって、\(\gcd(m,n)=1\) として

\[
\zeta=e^{-2\pi im/n}
\]

と書ける。

\subsection{B.4
チャネル振幅と二乗重み}\label{b4-ux30c1ux30e3ux30cdux30ebux632fux5e45ux3068ux4e8cux4e57ux91cdux307f}

\section{\$\$ U}\label{-u-1}

\section{\textbackslash frac12 \textbackslash begin\{pmatrix\}
1+\textbackslash zeta\&1-\textbackslash zeta\textbackslash{}
1-\textbackslash zeta\&1+\textbackslash zeta
\textbackslash end\{pmatrix\}}\label{frac12-beginpmatrix-1zeta1-zeta-1-zeta1zeta-endpmatrix-1}

\textbackslash begin\{pmatrix\} r\&t\textbackslash{} t\&r
\textbackslash end\{pmatrix\}, \$\$

ゆえに

\[
r=\frac{1+\zeta}{2},
\qquad
t=\frac{1-\zeta}{2}.
\]

その絶対値二乗は

\[
\boxed{
|r_{n,m}|^2
=\cos^2\left(\frac{\pi m}{n}\right)
}
\]

および

\[
\boxed{
|t_{n,m}|^2
=\sin^2\left(\frac{\pi m}{n}\right)
}
\]

である。これは、有限回で完全再帰する交換対称二チャネル過程の離散Born型重みである。

\subsection{B.5
数値実装座標の回収}\label{b5-ux6570ux5024ux5b9fux88c5ux5ea7ux6a19ux306eux56deux53ce}

実装で用いた

\[
t=e^{i\theta}\cos\theta,
\qquad
r=-ie^{i\theta}\sin\theta
\]

は

\[
\zeta=-e^{2i\theta}
\]

と置いた同じ作用素の座標表示である。実際、B.4の和・差は

\section{\$\$
\textbackslash frac\{1+\textbackslash zeta\}\{2\}}\label{-frac1zeta2}

\section{\textbackslash frac\{1-e\^{}\{2i\textbackslash theta\}\}\{2\}}\label{frac1-e2itheta2}

-ie\^{}\{i\textbackslash theta\}\textbackslash sin\textbackslash theta=r,
\$\$

\section{\$\$
\textbackslash frac\{1-\textbackslash zeta\}\{2\}}\label{-frac1-zeta2}

\section{\textbackslash frac\{1+e\^{}\{2i\textbackslash theta\}\}\{2\}}\label{frac1e2itheta2}

e\^{}\{i\textbackslash theta\}\textbackslash cos\textbackslash theta=t
\$\$

を与える。有限位数根では

\section{\$\$ \textbackslash theta}\label{-theta-1}

\textbackslash frac\{\textbackslash pi\}\{2\}-\textbackslash frac\{\textbackslash pi
m\}\{n\} \textbackslash pmod\textbackslash pi \$\$

となり、B.4の重みをそのまま再現する。

\begin{center}\rule{0.5\linewidth}{0.5pt}\end{center}

\section{付録C.
中心実装との対応}\label{ux4ed8ux9332c-ux4e2dux5fc3ux5b9fux88c5ux3068ux306eux5bfeux5fdc}

有限位数根をコードへ隠れた定数として入力していないこと、および線形核と後段写像の分離を確認できるよう、中心更新を示す。

\subsection{C.1 System
A：線形交換後のチャネル別正規化}\label{c1-system-aux7ddaux5f62ux4ea4ux63dbux5f8cux306eux30c1ux30e3ux30cdux30ebux5225ux6b63ux898fux5316}

\href{../20260715/run_system_A_localization_exchange_R_sweep_preliminary_v1.py}{System
A実装}の反復更新は次である。

\begin{Shaded}
\begin{Highlighting}[]
\NormalTok{a\_next }\OperatorTok{=}\NormalTok{ src.normalize(r }\OperatorTok{*}\NormalTok{ a }\OperatorTok{+}\NormalTok{ t }\OperatorTok{*}\NormalTok{ b)}
\NormalTok{b\_next }\OperatorTok{=}\NormalTok{ src.normalize(t }\OperatorTok{*}\NormalTok{ a }\OperatorTok{+}\NormalTok{ r }\OperatorTok{*}\NormalTok{ b)}
\NormalTok{a, b }\OperatorTok{=}\NormalTok{ a\_next, b\_next}
\end{Highlighting}
\end{Shaded}

ここには共通線形核

\[
(a,b)\mapsto(ra+tb,\ ta+rb)
\]

と、各出力チャネルへの個別正規化が明示されている。したがって局在性移乗の全写像は非線形である。

\subsection{C.2 System
B：無摂動有限位数核}\label{c2-system-bux7121ux6442ux52d5ux6709ux9650ux4f4dux6570ux6838}

\href{../20260715/run_minimal_system_B_gray_direct_check_v5.py}{System
B有限位数直接検査}の反復更新は次である。

\begin{Shaded}
\begin{Highlighting}[]
\NormalTok{a, b }\OperatorTok{=}\NormalTok{ normalize\_pair(r }\OperatorTok{*}\NormalTok{ a }\OperatorTok{+}\NormalTok{ t }\OperatorTok{*}\NormalTok{ b, t }\OperatorTok{*}\NormalTok{ a }\OperatorTok{+}\NormalTok{ r }\OperatorTok{*}\NormalTok{ b)}
\end{Highlighting}
\end{Shaded}

\(U_R\)
がユニタリで入力ペアが規格化されている限り、\texttt{normalize\_pair}
の規格化係数は1である。したがって、この直接検査でピーク位置を決める反復は厳密に線形核
\(U_R\) であり、\(R_{124,23}\) または \(R_{122,23}\)
を目標値として更新式へ注入していない。

\subsection{C.3 System
B：D強観測の非線形反作用}\label{c3-system-bdux5f37ux89b3ux6e2cux306eux975eux7ddaux5f62ux53cdux4f5cux7528}

\href{../20260714/run_gray_cat_d_observation_response_preliminary_v1.py}{D強観測実装}では、交換後に次の反作用を加える。

\begin{Shaded}
\begin{Highlighting}[]
\NormalTok{s\_next }\OperatorTok{=}\NormalTok{ s }\OperatorTok{+}\NormalTok{ d\_gain }\OperatorTok{*}\NormalTok{ s\_d }\OperatorTok{*}\NormalTok{ (}\FloatTok{1.0} \OperatorTok{{-}}\NormalTok{ s }\OperatorTok{*}\NormalTok{ s)}
\NormalTok{a\_next }\OperatorTok{=}\NormalTok{ math.sqrt(}\FloatTok{0.5} \OperatorTok{*}\NormalTok{ (}\FloatTok{1.0} \OperatorTok{+}\NormalTok{ s\_next)) }\OperatorTok{*}\NormalTok{ phase\_a}
\NormalTok{b\_next }\OperatorTok{=}\NormalTok{ math.sqrt(}\FloatTok{0.5} \OperatorTok{*}\NormalTok{ (}\FloatTok{1.0} \OperatorTok{{-}}\NormalTok{ s\_next)) }\OperatorTok{*}\NormalTok{ phase\_b}
\end{Highlighting}
\end{Shaded}

したがって、有限位数根は \(U_R\)
の円分スペクトルから生じ、A/Bへの強い選択は別のD反作用から生じる。コード構造そのものが、

\[
\boxed{
\text{再帰根の生成}
\ne
\text{強観測による状態選択}
}
\]

という本稿の二層分類を実装している。

\end{document}
